% =====================================================================
% 1. BUSINESS UNDERSTANDING
% Rubric: problem/goal from a business perspective, purpose, objectives.
% Page budget: \todo{decide}
% =====================================================================
\section{Business Understanding}
\label{sec:business}
The report follows the CRISP-DM process \citep{chapman2000crisp}.

% --- 1.1 Context -----------------------------------------------------
% Why a bank cares about EPC labels: ESG reporting, green mortgages,
% climate/transition risk in the collateral portfolio, regulation.
% Cite sources (EU taxonomy / EBA / EPBD) in references.bib.
\paragraph{Context.}
\todo{Why the EPC label of a property matters to a bank (ESG, risk, products).}

% --- 1.2 Problem -----------------------------------------------------
% The label is not always available -> needs to be estimated.
\paragraph{Problem.}
\todo{What is missing today and what decision it blocks.}

% --- 1.3 Data-mining goal --------------------------------------------
% Translate the business goal into a ML task.
\paragraph{Data-mining goal.}
Predict the \epc{} label on the ordered scale
A+ $>$ A $>$ B $>$ C $>$ D $>$ E $>$ F $>$ G (eight ordinal classes)
for every record of the scoring set.
\todo{One sentence on why this is ordinal multiclass, not plain classification.}

% --- 1.4 Success criteria --------------------------------------------
\paragraph{Success criteria.}
The official metric is \onenotch{}: a prediction is correct when it is
within one class of the true label, i.e.\
$\text{acc}_{1} = \frac{1}{n}\sum_i \mathbb{1}\big[|r(\hat y_i) - r(y_i)| \le 1\big]$,
with $r(\cdot)$ the rank A+\,=\,0, \dots, G\,=\,7.
Exact accuracy and the mean absolute rank error are reported as secondary metrics.
\todo{Business reading: why a 1-class error is acceptable and a 2+-class error is not.}

% --- 1.5 Objectives / research questions -----------------------------
\paragraph{Objectives.}
\begin{enumerate}
  \item \todo{O1: accurate predictions (metric above).}
  \item \todo{O2: which groups of information drive the label (explainability).}
  \item \todo{O3: actionable recommendation for the bank.}
\end{enumerate}
